\newpage
\section*{第23章：全称类型}
\addcontentsline{toc}{section}{第23章：全称类型}

\begin{taplsolutionentrybox}
\phantomsection\label{sol:translator:23.4.1}
\textbf{23.4.1 解答}

以 \(\taplterm{selfApp}\) 为例。上下文中
\[
x:\forall X.\,X\to X.
\]
由 T-TApp，在类型 \(\forall X.\,X\to X\) 处实例化 \(x\)，得到
\[
x[\forall X.\,X\to X]:
(\forall X.\,X\to X)\to(\forall X.\,X\to X).
\]
再把原来的 \(x\) 作为实参应用，即得结果类型
\(\forall X.\,X\to X\)。外层 T-Abs 给出题中函数类型。
\(\taplterm{id}\)、\(\taplterm{double}\) 与
\(\taplterm{quadruple}\) 逐层使用 T-TAbs、T-TApp、T-Abs 和 T-App 即可。
\taplsolutionbacklink{ex:23.4.1}{返回练习23.4.1}
\end{taplsolutionentrybox}

\begin{taplsolutionentrybox}
\phantomsection\label{sol:translator:23.4.2}
\textbf{23.4.2 解答}

固定 \(X,Y\) 后，\(f:X\to Y\)、\(l:\tapltype{List}\,X\)。
空分支 \(\taplterm{nil}[Y]\) 具有类型 \(\tapltype{List}\,Y\)。
非空分支中
\[
\taplterm{head}[X]l:X,\qquad
f(\taplterm{head}[X]l):Y,
\]
递归调用
\[
m(\taplterm{tail}[X]l):\tapltype{List}\,Y.
\]
因此 \(\taplterm{cons}[Y]\) 的两个实参类型正确，两个条件分支都具有
\(\tapltype{List}\,Y\)。fix 的函数体由此具有
\(\tapltype{List}\,X\to\tapltype{List}\,Y\)，再抽象 \(f,Y,X\) 即得
所示全称类型。
\taplsolutionbacklink{ex:23.4.2}{返回练习23.4.2}
\end{taplsolutionentrybox}

\begin{taplsolutionentrybox}
\phantomsection\label{sol:translator:23.4.4}
\textbf{23.4.4 解答}

可采用插入排序：先定义
\[
\taplterm{insert}:
\forall X.\,(X\to X\to\tapltype{Bool})\to
X\to\tapltype{List}\,X\to\tapltype{List}\,X,
\]
它沿列表递归，在第一个满足比较条件的位置插入元素；再以空列表为初值，把
原列表的每个元素依次插入。若允许 \(\taplterm{fix}\)，两个递归都可直接按
\cref{sec:23.4}中 \(\taplterm{map}\) 的形式书写；若要求纯 System F，
则使用\cref{sol:author:23.4.12}的折叠编码。
\taplsolutionbacklink{ex:23.4.4}{返回练习23.4.4}
\end{taplsolutionentrybox}

\begin{taplsolutionentrybox}
\phantomsection\label{sol:translator:23.4.7}
\textbf{23.4.7 解答}

对 \(\taplterm{ctimes}\)，固定 \(X\) 和 \(s:X\to X\) 后，
\(m[X]s:X\to X\)。Church 数 \(n\) 把这个函数迭代 \(n\) 次；每次又会把
\(s\) 迭代 \(m\) 次，所以总计迭代 \(mn\) 次。

对 \(\taplterm{cexp}\)，\(m[X]: (X\to X)\to X\to X\)，因此可把它看成
类型 \(X\to X\) 上的“迭代器”。令 \(n\) 在 \(X\to X\) 这一类型上迭代
\(m[X]\)，便把“乘以 \(m\) 次迭代次数”的操作再迭代 \(n\) 次，得到
\(m^n\) 次应用。逐层使用 T-TApp 可得两项均返回 \(\tapltype{CNat}\)。
\taplsolutionbacklink{ex:23.4.7}{返回练习23.4.7}
\end{taplsolutionentrybox}
